%% Define document class, variables and frontmatter
\documentclass{beamer}
\title{Progress report}
%%% Variables, functions and other settings
%% Define beamer theme
\usetheme{Goettingen}
%% Packages
%\usepackage[utf8]{inputenc}
%\usepackage[T1]{fontenc}
%% Define basic variables
\subtitle{}
\author{Neronet}
\institute[]{
  \emph{
    Toolbox for managing the training \\
    neural networks
  } \\[0.5cm]
  CSE-C2610 \\
  Software Project \\[0.2cm]
  Aalto University
}
\date{\today{}}
\subject{Software engineering}
%% Colors and visuals
% Specify link and URL colors
\hypersetup{colorlinks,linkcolor=,urlcolor=magenta}
% Replace navigation symbols with a logo
\logo{\includegraphics[height=1cm]{gfx/logo_aalto.png}}
\setbeamertemplate{navigation symbols}{\insertlogo}
%% Functions
% Define short set and clear background commands
\newcommand{\bgset}[1]{\usebackgroundtemplate{
  \includegraphics[width=\paperwidth,height=\paperheight]{#1}}}
\newcommand{\bgclear}{\usebackgroundtemplate{}}
\ifdefined \showoutlineatsection
  % Have LaTeX render an outline frame just before each new section
  \AtBeginSection[]{
    \bgset{gfx/neural3__bgmod.jpg}
    \begin{frame}<beamer>{Outline}
      \tableofcontents[currentsection]
    \end{frame}
    \bgset{gfx/neural1__bgmod.jpg}}
\fi
%%% Main document content
\begin{document}
%% Title page
\bgset{gfx/neural2__bgmod.jpg}
\begin{frame}
  \titlepage
\end{frame}
%% Outline page
\bgset{gfx/neural3__bgmod.jpg}
\begin{frame}{Outline}
  \tableofcontents
\end{frame}
\def\Done{\textcolor{green}{Done}}
%% Content pages
\section{Introduction} \bgset{gfx/neural1__bgmod.jpg}
\begin{frame}{Introduction}{Goals} % Joona
  Our goal is to develop a tool for computational researchers to enable easy
  \begin{itemize}
  \item specification and management of experiment queues
  \pause \item batch submission of experiment jobs to computing clusters
  \pause \item monitoring of ongoing experiments' logs and parameter values
  \pause \item access to experiment information during and after the run
  \pause \item configurable notifications on experiment state and progress
  \pause \item configurable criteria for experiment autotermination
  \pause \item logging of experiment history
  \end{itemize}
\end{frame}
\begin{frame}{Introduction}{What} % Samuel
  In essence the product is a Python-based tool that enables computational researchers to conduct their research more effectively.
  \begin{itemize}
  \item \pause It utilizes SSH and TCP sockets to distribute the computational workload into computer clusters. It supports the Slurm job and resource manager but can function without it as well.
  \item \pause It is framework agnostic in that it permits the use of a very wide variety of tools to actually conduct the computing.
  \end{itemize}
\end{frame}
\section{Results} % Samuel & Joona
\begin{frame}{Results}{Sprint 0} \bgset{gfx/neural3__bgmod.jpg}
  Goal: \emph{Team building and preparing for sprint 1} \pause \Done{}

  \pause Product Backlog Items: \emph{None}

  \pause Sprint 0 took a lot of effort from us since the project topic was very challenging to dive into. Also none of us had done this course before. Interviews with Jelena \& Simo helped us to understand the project.
  
  \pause We were proud of our efforts in the sprint.
\end{frame}
\begin{frame}{Results}{Sprint 1}
  Goal: \emph{Develop a prototype that offers the most basic functionality via a CLI}
  
  \pause Product Backlog Items: 
  \begin{itemize}
  \pause \item US1: As a user, I want a basic user guide that would cover the installation of Neronet and its use via CLI. \pause \Done{}
  \pause \item US2: As a user, I want to specify clusters by address and type to specify my computing resources. \pause \Done{}
  \pause \item US3: As a user, I want to specify experiments by name, files and parameters. \pause \Done{}
  \pause \item US4: As a user, I want to submit experiments to unmanaged nodes. \pause \Done
  \pause \item US5: As a user, I want an experiment status report so that I can review experiment status details. \pause \Done{}
  \end{itemize}

  \pause Just a prototype, not ready for release to users.
\end{frame}
\section{Demo}
\begin{frame}{Demo} % Tuomo & Teemu
  Demo script:
  \begin{enumerate}
  \item Neronet Installation, preferences and initial setup of clusters
  \item Specification of clusters via CLI
  \item Specification of an experiment
  \item Submission of the specified experiment to an unmanaged node
  \item Retrieval of experiment status report
  \end{enumerate}
\end{frame}
\section{Quality} \bgset{gfx/neural1__bgmod.jpg}
\begin{frame}{Quality} % Matias & Iiro
  Definition of done:
  \begin{itemize}
  \item We defined \Done{} in three levels: BI, sprint and project
  \item BI level: unit tests done where applicable,
    functional test coverage 80\%, conformity (PEP-8), commented,
    documented, peer reviewed
  \item Sprint level: BI:s are \Done{}, increment is tested and reviewed,
    sprint goal is achieved
  \pause \item Updates to DoD:
    \begin{itemize}
    \item We replaced \emph{unit test coverage 90\%} with \emph{unit tests
      are written where applicable} -- the old metric was not useful for all
      BIs
    \end{itemize}
  \item Otherwise, we have followed our DoD as planned.
  \end{itemize}
\end{frame}
\begin{frame}{Quality} % Matias & Iiro
  Used QA practices and tools:
  \begin{itemize}
  \item Commenting \& documentation -- forces to rethink from another
    perspective, facilitates peer review
  \item Python standard unittest framework -- white-box test automation
  \item Functional testing -- manual black-box testing based on specs
  \item Peer review -- quality assurance
  \end{itemize}
\end{frame}
\begin{frame}{Quality} % Matias & Iiro
  Performance in quality assurance practices peer-graded with scale 0-3:
  \begin{table} \small
  \begin{tabular}{l|rrrrrrrrrrr}
    US &   UT &   FT &  Com &  Doc & Peer  \\ \hline
     1 &    3 &    3 &    3 &    3 &    3  \\
     2 &    2 &    2 &    3 &    2 &    3  \\
     3 &    3 &    3 &    3 &    3 &    3  \\
     4 &    3 &    3 &    3 &    3 &    1  \\
     5 &    1 &    3 &    2 &    1 &    1  \\ \hline
       &    2 &    3 &    3 &    2 &    2
  \end{tabular}
  \end{table}
  \pause Qualitatively we achieved our standards only partially:
  \begin{itemize}
  \item Unit and functional test coverage -- good
  \item Quality of comments and documentation -- good
  \item Peer review -- ok (done rather quickly)
  \end{itemize}
\end{frame}
\begin{frame}{Quality} % Matias & Iiro
  Relevant quality attributes:
  \begin{itemize}
  \pause \item \alert{Usability} -- Our user guide helps newbies understand
    our product $\rightarrow$ The usability should be good
  \pause \item \alert{Reliability} -- We estimate the overall reliability
    only satisfactory. We will put more effort on this in the future.
  \pause \item \alert{Extendability} -- At the moment our software's
    extendability is ok, but it is only a prototype.
  \pause \item \alert{Performance} -- The current performance is mediocre.
    We will put more effort on this in the future.
  \end{itemize}
\end{frame}
\section{Effort}
\begin{frame}{Effort} % Matias & Iiro
Spent and budgeted effort in hours by team member and sprint:
\begin{table} \small
  \begin{tabular}{l|rrrrrrrrrrr}
    S &  Samuel &   Teemu &   Tuomo &   Joona &    Iiro &  Matias  \\ \hline
    0 &  140/50 &   36/35 &   45/35 &   40/35 &   36/35 &   43/35  \\
    1 &   52/30 &   37/33 &   42/33 &   46/33 &   32/33 &   37/33  \\
    2 &    0/30 &    0/33 &    0/33 &    0/33 &    0/33 &    0/33  \\
    3 &    0/15 &    0/33 &    0/33 &    0/33 &    0/33 &    0/33  \\
    4 &    0/15 &    0/33 &    0/33 &    0/33 &    0/33 &    0/33  \\
    5 &    0/15 &    0/33 &    0/33 &    0/33 &    0/33 &    0/33  \\
    6 &    0/20 &    0/25 &    0/25 &    0/25 &    0/25 &    0/25  \\ \hline
      & 192/175 &  73/225 &  87/225 &  86/225 &  68/225 &  80/225  \\
  \end{tabular}
  \end{table}
  \pause \begin{itemize}
  \item Our team had comparatively little experience
  \item The project's problem domain was challenging to understand
  \item Studying existing products and technologies took a lot of time
  \item The scrum master served also as team leader and lead developer
  \end{itemize}
\end{frame}
\section{Retros} \bgset{gfx/neural3__bgmod.jpg}
\begin{frame}{Retros: Sprint 0} % Joona & Samuel
  Sprint planning:
  \begin{itemize}
  \item BI clarity and simplicity (user guide helps)
  \item It might have been better if the PO had created the stories from scratch -Matias, Tuomo
  \item We should actively seek more input from PO when developing the user guide
  \item we should make sure we reserve enough time for the actual story selection on Monday -Matias
  \end{itemize}
\end{frame}
\begin{frame}{Retros: Sprint 0} % Joona & Samuel
  Daily scrums:
  \begin{itemize}
  \item We have mostly been doing teamwork, so there has been little new info in the scrums -Matias -Joona -Teemu
  \item They have been overly long and they have extended due to inexperience.
  \item People are late.
  \end{itemize}
\end{frame}
\begin{frame}{Retros: Sprint 0} % Joona & Samuel
  Teamwork sessions:
  \begin{itemize}
  \item Sessions are too long and sometimes people get hungry
  \item Generally someone has to leave early or comes late
  \end{itemize}
\end{frame}
\begin{frame}{Retros: Sprint 0} % Joona & Samuel
  Tools:
  \begin{itemize}
  \item Flowdock is good x6
  \item Agilefant has a steep learning curve. -Iiro
  \item People tend to forget to log their time at agilefant
  \item Github hasn't been used much. Hope to use it more during future sprints
  \item Floobits is very buggy.
  \item Top 3 tools: 1) GitHub 2) Flowdock 3) Agilefant
  \item Worst 3 tools: 1) Floobits 2) Six tactics 3) Agilefant
  \end{itemize}
\end{frame}
\begin{frame}{Retros: Sprint 0} % Joona & Samuel
  How teamwork could be improved:
  \begin{itemize}
  \item People should be more on time
  \item Scrum Master shouldn't have to work as a team leader too.
  \item Hard to think of improvements since we haven't really started coding yet
  \end{itemize}
\end{frame}
\end{document}